%% ============================================================
%% 分片 24：原文第 9 节「't Hooft And Hecke Operators」（原文 glshort.tex 起于第 6755 行）
%% 由 tools/harvmac2tex.py 从 arXiv 源码（hep-th/0604151v3）自动生成。
%% 只译英文散文；公式、编号、\label、\cite 一律不改。
%% ============================================================

\subsection{Hecke 修正的空间}

\label{sec:spaceheck}

本节余下部分致力于更仔细地考察这个问题中的全纯数据。

我们将先给出全纯数据的直接描述，着重于简单的例子，然后在第~\ref{sec:affgrass}~节 解释数学文献中通常用来安置 Hecke 修正空间的更有力的框架。

这些讨论都是局部的，因此为简单起见，我们取 $C=\Bbb{CP}^1$，并将它视为复 $z$-平面再加上无穷远点。我们取点 $p_0$ 为 $z=0$。首先，取规范群为 $G=U(N)$。设 $E_-$ 是秩为 $n$ 的平凡丛。（由于讨论是局部的，初始取哪个 $E_-$ 并无关系。）我们要讨论 't Hooft/Hecke 算子如何修改 $E_-$，从而得到一个新的丛 $E_+$。我们把 $E_-$ 的一个局部全纯截面（即 $p_0$ 附近全纯的截面）写成 $s=(f_1,f_2,\dots,f_N)$，其中 $f_i$ 是局部全纯函数。

\begin{yizhu}[原文疑误]
按：原文写「a trivial bundle of rank $n$」（秩为 $n$ 的平凡丛），但全篇随后一律用 $N$ 记这个秩：$E_-=\oplus_{i=1}^N{\cal L}_i$、$s=(f_1,\dots,f_N)$、$G=U(N)$。此处应为 $n=N$；译文照录 $n$，读者按秩 $N$ 理解即可。
\end{yizhu}

若 $s$ 是 $E_-$ 在 $p_0$ 附近的任一全纯截面，我们可以如下在 $p_0$ 附近定义一个线丛 ${\cal L}$。${\cal L}$ 是 $E_-$ 的子丛，它在 $p_0$ 的某个邻域上由下述方式定义：${\cal L}$ 的截面就是 $E_-$ 的、形如 $gs$ 的截面，其中 $g$ 是任意全纯函数。我们称 ${\cal L}$ 为由 $s$ 生成的线丛。现在设 $s_1,\dots,s_N$ 是 $E_-$ 的 $N$ 个局部全纯截面，它们在 $z=0$ 上的限制线性无关。（若记 $s_i(z)=(f_{1i}(z),f_{2i}(z),\dots,f_{Ni}(z))$，条件是 $N\times N$ 矩阵 $f_{ji}(0)$ 的行列式非零。）对每个 $i$，令 ${\cal L}_i$ 为由 $s_i$ 生成的线丛。线性无关性保证在 $z=0$ 附近有分解 $E_-=\oplus_{i=1}^N{\cal L}_i$，这等价于说 $E_-$ 的任一局部全纯截面 $s$ 都可写成 $s=\sum_i g_is_i$，其中 $g_i$ 是某些局部全纯函数。在这种情形下，我们说截面 $s_i$ 生成 $E_-$。

反过来，$E_-$ 的任何这样的局部分解都取上述形式。若 $E_-=\oplus_{i=1}^N{\cal L}_i$，则对每个 $i$ 取 ${\cal L}_i$ 的任一在 $z=0$ 处非零的截面 $s_i$。于是在 $z=0$ 附近，${\cal L}_i$ 就是由 $s_i$ 生成的线丛，从而 $E_-$ 的分解取上一段所描述的形式。

现在考虑 't Hooft/Hecke 算子 $T(^L\negthinspace w)$ 的作用，其中 $^L\negthinspace w=(m_1,m_2,\dots,m_N)$。对于我们所选定的 $E$ 的分解，它把 $E_-$ 映到这样一个丛：在 $z=0$ 附近为 $E_+=\oplus_i\left( {\cal L}_i\otimes {\cal O}(p_0)^{m_i}\right)$。在 $z=0$ 之外，$E_+$ 与 $E_-$ 完全相同。

但 ${\cal L}_i\otimes {\cal O}(p_0)^{m_i}$ 是这样一条线丛：它在 $z=0$ 附近的一般截面为 $s_i/z^{m_i}$。于是 $E_+$ 在 $z=0$ 附近的一般截面取形式 $\sum_{i=1}^N g_i(z) s_i(z)z^{-m_i}$，其中函数 $g_i$ 在 $z=0$ 附近全纯。在 $z=0$ 之外，$E_+$ 的全纯截面与 $E_-$ 的全纯截面相同：它具有形式 $(f_1,\dots,f_N)$，其中 $f_1,\dots,f_N$ 是任意全纯函数。

不同的分解可能导致同一个 $E_+$。这正是我们要研究的问题：我们想研究能被给定的 't Hooft 算子把 $E_-$ 映过去的所有丛 $E_+$ 构成的空间。在表述这个问题时，我们必须记住：$E_-$ 的任一 Hecke 修正 $E_+$ 都在 $z=0$ 之外自然同构于 $E_-$，所以若 $E_+$ 与 $\tilde E_+$ 是 $E_-$ 的两个 Hecke 修正，它们在 $z=0$ 之外彼此自然同构。我们称 $E_+$ 与 $\tilde E_+$ 等价，当且仅当 $z=0$ 之外的这个自然同构可以延拓过 $z=0$。我们要把可能的 $E_+$ 按此等价关系分类。这种分类在实践中意味着什么，将通过考察例子而变得清楚。

\begin{yizhu}[易错点]
按：这里分类的是「丛 $E_+$ 连同它在 $z=0$ 之外到 $E_-$ 的自然同构」，而不是单独的丛同构类。两个 Hecke 修正只有在那个自然同构能延拓过 $z=0$ 时才等价；抽象地同构但延拓不过去的不算同一个点。后文诸例中出现的各种等价关系（整体缩放、$u'$ 方向的平移等）都是这一条的具体表现。
\end{yizhu}

我们从 $N=2$ 的几个简单例子开始。首先取 $^L\negthinspace w=(1,0)$。设 $s_1(z)=(e_1(z),f_1(z))$、$s_2(z)=(e_2(z),f_2(z))$，且 $(e_1(0),f_1(0))$ 与 $(e_2(0),f_2(0))$ 线性无关。$E_+$ 由 $s_1'=z^{-1}s_1$ 连同 $s_2$ 生成。特别地，$s_1=zs_1'$ 是 $E_+$ 的截面，$s_2$ 当然也是。由于 $s_1$ 与 $s_2$ 生成 $E_-$，这意味着 $E_-$ 的任一截面都是 $E_+$ 的截面。但 $E_+$ 的截面可以有正比于 $z^{-1}s_1(0)$ 的极部。

总之，在这个例子中，$E_+$ 的任一截面都取形式
% 式 (9.16)
\setcounter{equation}{15}
\begin{equation}\label{eq:milok}
s(z)={c\over z}s_1(0)+(h_1(z),h_2(z))
\end{equation}
在 $z=0$ 附近，其中 $c$ 是复常数，$h_1$、$h_2$ 是局部全纯函数。显然，$E_+$ 只依赖于 $s_1(0)=(e_1(0),f_1(0))$ 的选取，且相差一个复缩放。所以我们应把 $(e_1(0),f_1(0))$ 看作定义了 $\Bbb{CP}^1$ 中的一个点，而这个 $\Bbb{CP}^1$ 参数化了所有能由 $E_-$ 经 $T(^L\negthinspace w)$（$^L\negthinspace w=(1,0)$）作用构造出的丛 $E_+$。

我们将说 $\Bbb{CP}^1$ 是型为 $^L\negthinspace w=(1,0)$ 的所有 Hecke 修正构成的空间。把这个空间记作 ${{\EUY}}(^L\negthinspace w)$，或在需要指明作 Hecke 修正的点（此处为 $p_0$）时记作 ${{\EUY}}(^L\negthinspace w,p_0)$。我们也把同一个空间记作 ${{\EUY}}(^L\negthinspace R)$，其中 $^L\negthinspace R$ 是 $^L\negthinspace G$ 的以 $^L\negthinspace w$ 为最高权的表示，并又称之为型为 $^L\negthinspace R$ 的 Hecke 修正空间。

在这个具体的例子里，Hecke 修正空间是紧的且光滑的。但这并不典型。再看另一个例子：取 $^L\negthinspace w=(2,0)$。Hecke 算子 $T(^L\negthinspace w)$ 作用在与前面相同分解下的 $E_-$ 上，现在产生由 $s_1'=z^{-2}s_1$ 与 $s_2$ 生成的丛 $E_+$。由于 $s_1=z^2s_1'$ 与 $s_2$ 一样是 $E_+$ 的截面，$E_-$ 的任一截面 $(h_1,h_2)$ 仍然定义 $E_+$ 的一个截面。此外，$E_+$ 还有一个二维的极截面空间，由 $z^{-2}s_1$ 与 $z^{-1}s_1$ 生成。记 $e_1(z)=u+zu'+{\cal O}(z^2)$、$f_1(z)=v+zv'+{\cal O}(z^2)$，其中 $u,u',v,v'\in\Bbb{C}$。则 $E_+$ 的一般局部全纯截面取形式 % 式 (9.17)
\setcounter{equation}{16}
\begin{equation}\label{eq:ilko}
s(z)={c\over
z^2}(u+zu',v+zv')+{d\over z} (u,v)+(h_1(z),h_2(z)),
\end{equation}
其中 $c$ 与 $d$ 是复常数，$h_1(z),h_2(z)$ 是全纯函数。

如何描述这个情形下的 Hecke 修正空间？显然，对这四个变量作整体缩放
% 式 (9.18)
\setcounter{equation}{17}
\begin{equation}\label{eq:lorky}
(u,u',v,v')\to\lambda (u,u',v,v')
\end{equation}
是无关紧要的；它可以吸收进 $c$ 与 $d$ 的缩放中。但还有另一个我们应当考虑的等价关系。若我们作变换 % 式 (9.19)
\setcounter{equation}{18}
\begin{equation}\label{eq:morky}
(u',v')\to (u', v')+w(u,v),
\end{equation}
其中 $w\in\Bbb{C}$，那么这可以吸收进 $d\to d-wc$。所以我们应把这个变换看作参数 $u,u',v,v'$ 上的另一个等价关系。这两种对称性都源于这样的事实：$E_+$ 在 $s_1\to gs_1$ 下不变，其中 $g$ 是在 $z=0$ 处可逆的全纯函数。


于是 Hecke 修正空间由 $(u,u',v,v')$ 参数化，其中 $u$ 与 $v$ 不能同时为零（因为 $s_1(0)$ 必须非零），并且要受等价关系 \eqref{eq:lorky}\ 与 \eqref{eq:morky} 的约束。在变换 \eqref{eq:morky}\ 下的不变量是 $u$、$v$ 与 $b=vu'-uv'$。因此，我们可以直接用 $u$、$v$、$b$ 来参数化 Hecke 修正，从而把这个缩放考虑进去。$b$ 在 \eqref{eq:lorky}\ 下按权 2 缩放：$b\to \lambda^2b$。于是任何不全为零的三元组 $(u,v,b)$ 在模去这个缩放后，定义了加权射影空间 $\Bbb{WCP}^2(1,1,2)$ 中的一个点。我们迄今定义的 Hecke 修正空间并不是整个加权射影空间：我们漏掉了单独一个点 $(0,0,1)$，因为我们不允许 $u=v=0$。

尽管可以在非紧的 Hecke 修正空间上工作，但这并不方便。$A$-模型中的物理态空间是各种模空间的上同调群。紧性使相应上同调群的定义直接得多；否则就需要详细讨论波函数在无穷远处（在目前情形即 $(0,0,1)$ 附近）应当如何行为。

我们迄今定义的 Hecke 修正空间 ${{\EUY}}(^L\negthinspace w)$ 确实有一个自然的紧化，它来自这样一个事实：当 $u,v\to 0$ 时丛 $E_+$ 有自然的极限。引入小参数 $\epsilon$，令 $u=\epsilon \bar u$、$v=\epsilon \bar v$、$d=d'/\epsilon$。$E_+$ 的一般截面取形式 % 式 (9.20)
\setcounter{equation}{19}
\begin{equation}\label{eq:gilko}
s(z)={c\over
z^2}(\epsilon \bar u+zu',\epsilon \bar v+zv')+{d'\over z} (\bar
u,\bar v)+(h_1(z),h_2(z)).
\end{equation}
令 $\epsilon\to 0$ 而保持其他变量不变，其极限正是 % 式 (9.21)
\setcounter{equation}{20}
\begin{equation}\label{eq:pilko}
s(z)={c\over z}(u',v')+{d'\over z}(\bar u,\bar
v)+(h_1(z),h_2(z)).
\end{equation}

一般截面取这种形式的丛可以有更简单的解释。加权射影空间上的点 $(0,0,1)$ 满足 $b\not=0$，这等价于向量 $(u,v)$ 与 $(u',v')$ 线性无关，也等价于 $(\bar u,\bar v)$ 与 $(u',v')$ 线性无关。所以当我们通过 $\epsilon\to 0$ 的极限逼近 $(0,0,1)$ 时，应当假定这种线性无关性。既然如此，\eqref{eq:pilko}\ 可以更简单地描述为：$s(z)$ 在 $z=0$ 处可以有任意的一阶极点。从而，在这个极限下的丛就是 $E_+=E_-\otimes {\cal O}(p_0)$。

但这正是我们用 $^L\negthinspace w=(1,1)$ 对丛 $E_-$ 作 Hecke 修正时所得的结果。对这个 $^L\negthinspace w$，Hecke 运算 $T(^L\negthinspace w)$ 就是与 ${\cal O}(p_0)$ 作张量积的运算，这已在第~\ref{sec:theck}~节 说明过。于是结论是：$^L\negthinspace w=(2,0)$ 的 Hecke 修正空间有一个自然的紧化，其中也允许 $^L\negthinspace w=(1,1)$ 的 Hecke 修正。对于型为 ${{\EUY}}(^L\negthinspace w)$ 的 Hecke 修正空间的这类紧化，我们将记作 $\bar {{\EUY}}(^L\negthinspace w)$（若希望指明作 Hecke 修正的点，则更明确地写作 $\bar {{\EUY}}(^L\negthinspace w,p_0)$）。

\begin{yizhu}[物理图像]
按：紧化的必要性来自 $A$-模型：其物理态空间取在各模空间的上同调上，而 $(0,0,1)$ 这个点对应 $u=v=0$（即 $s_1(0)=0$），Hecke 修正空间在它附近不再紧，态空间的上同调也就无从定义。补进来的极限正是 $E_+=E_-\otimes{\cal O}(p_0)$，即 $^L\negthinspace w=(1,1)$ 的 Hecke 修正。几何上，${\cal O}(2)\oplus{\cal O}$ 不稳定、可无穷小扰动到 ${\cal O}(1)\oplus{\cal O}(1)$，两者第一 Chern 类同为 $2$，所以补进来的点与原来的点属于同一拓扑型——这正是后文「不稳定的 't Hooft 算子可以混合」的几何根源。
\end{yizhu}

我们回忆，一般地，一个 't Hooft 算子由 $\Bbb{CP}^1$ 上一个 $G$-丛的选取来指定；给定这样的选取后，就要解 Bogomolny 方程，并要求有由该 $G$-丛的选取所确定的局部奇点。特别地，对 $G=U(2)$，具有 $^L\negthinspace w=(m_1,m_2)$ 的 't Hooft 算子是用 $\Bbb{CP}^1$ 上的 $U(2)$ 丛 ${\cal O}(m_1)\oplus {\cal O}(m_2)$ 来定义的。这种丛唯一的拓扑不变量是第一 Chern 类，它等于 $m_1+m_2$。对应于 $^L\negthinspace w=(2,0)$ 的丛 ${\cal O}(2)\oplus {\cal O}$ 是一个「不稳定丛」；它可以经无穷小扰动而把全纯型变为 ${\cal O}(1)\oplus {\cal O}(1)$，后者对应于 $^L\negthinspace w =(1,1)$。我们刚才描述的分析表明，为了紧化 $^L\negthinspace w=(2,0)$ 的 Hecke 修正空间，必须把对应于 $^L\negthinspace w=(1,1)$ 的 Hecke 修正也包括进来。

我们考虑过的另一个例子则相当不同。丛 ${\cal O}(1)\oplus {\cal O}(0)$ 不能扰动成任何别的全纯型。所以 $^L\negthinspace w=(1,0)$ 的 Hecke 修正空间是光滑且紧的。

\par\medskip{\it 加权射影空间的奇点}

为了更深入地理解 $^L\negthinspace w=(2,0)$ 的例子，让我们重新考察加权射影空间 $\Bbb{WCP}^2(1,1,2)$，它有权为 $(1,1,2)$ 的齐次坐标 $(u,v,b)$。在点 $(0,0,1)$ 处，这个空间有 $\Bbb{Z}_2$ 轨道流形奇点。

因此局部结构就是 $A_1$ 奇点 $\Bbb{C}^2/\Bbb{Z}_2$ 的结构。$\Bbb{C}^2/\Bbb{Z}_2$ 是超凯勒轨道流形。与典型的复奇点不同，$A_1$ 奇点可以被相容地同时解奇和形变。最自然的理解方式 \cite{ref121} 是把 $\Bbb{C}^2/\Bbb{Z}_2$ 实现为 $\Bbb{H}^2\cong \Bbb{C}^4$ 由 $U(1)$ 作用的超凯勒商。这里我们不走这条路，而只是手算描述加权射影空间 $\Bbb{WCP}^{2}(1,1,2)$ 奇点的形变与解奇。

为描述形变，令 $a_1=u^2,$ $a_2=uv$、$a_3=v^2$。变量 $(a_1,a_2,a_3,b)$ 都有权 2，模去缩放后它们参数化一个普通的射影空间 $\Bbb{CP}^3$。但它们满足 $a_1a_3-a_2^2=0$，所以 $\Bbb{WCP}^2(1,1,2)$ 就是由这个方程定义的 $\Bbb{CP}^3$ 的子簇。方程 $a_1a_3-a_2^2=0$ 可以形变为 $a_1a_3-a_2^2=\epsilon b^2$，其中 $\epsilon$ 是小参数。最后这个方程定义了 $\Bbb{CP}^3$ 中的一张光滑二次曲面，它同构于\footnote{通过线性变量替换，可以把 $a_1,a_2,a_3$ 与 $b$ 安排为 $2\times 2$ 矩阵 $M$ 的矩阵元，使方程变为 $\det(M)=0$。这个方程表示 $M$ 的秩为 1，因此具有形式 $M_{ij}=x_iy_j$，其中 $x_i$ 与 $y_j$ 分别是 $\Bbb{CP}^1\times \Bbb{CP}^1$ 中两个因子的齐次坐标。} $\Bbb{CP}^1\times \Bbb{CP}^1$，于是 $\Bbb{WCP}^2(1,1,2)$ 可以形变为 $\Bbb{CP}^1\times \Bbb{CP}^1$。

为描述解奇，我们从坐标为 $u,v,b,b'$ 的 $\Bbb{C}^4$ 出发。考虑 $U(1)\times U(1)$ 的作用 $(u,v,b,b')\to (\lambda u,\lambda v,\lambda^2 \tilde \lambda b,\tilde\lambda b')$，其中 $|\lambda|=|\tilde\lambda|=1$。定义 moment map $\mu=|u|^2+|v|^2+2|b|^2$、$\tilde\mu= |b|^2+|b'|^2$，并对 $\Bbb{C}^4$ 取 $U(1)^2$ 的辛商。为此，施加 moment map 方程
% 式 (9.22)
\setcounter{equation}{21}
\begin{equation}\label{eq:ripo}
\begin{aligned}|u|^2+|v|^2+2|b|^2 & = 1\\
                                        |b|^2+|b'|^2& = d\\\end{aligned}
\end{equation}

其中 $d$ 是某个常数，然后除以 $U(1)\times U(1)$。对小的正 $d$，这些方程不允许 $u=v=0$，所以可以把 $(u,v)$ 当作某个 $B=\Bbb{CP}^1$ 的齐次坐标。对小的 $d$，$\Bbb{CP}^1$ 上给定点处的纤维是另一个 $\Bbb{CP}^1$，由 $b$ 与 $b'$ 参数化。由于在 $(u,v)\to (\lambda u,\lambda v)$ 下 $(b,b')$ 变为 $(\lambda^2b,b')$，$b$ 与 $b'$ 分别在 $B$ 上取值于线丛 ${\cal O}(2)$ 与 ${\cal O}$。所以对小的正 $d$，辛商以 $B=\Bbb{CP}^1$ 为底纤维化，纤维为 $Y=\Bbb{P}({\cal O}(2)\oplus {\cal O})$，即秩二向量丛 $\CO(2)\oplus \CO$ 的射影化。特别地，对小的正 $d$，辛商是 $\Bbb{CP}^1$ 上的 $\Bbb{CP}^1$ 丛，因此是一张 Hirzebruch 曲面。

若增大 $d$，在到达 $d=1/2$ 之前什么也不会发生。在 $d=1/2$ 处，Hirzebruch 曲面出现奇点并退化为加权射影空间 $\Bbb{WCP}^2(1,1,2)$。在 $d=1/2$ 时，我们只需把第二个 moment map 方程写成 $|b'|^2=1/2-|b|^2$，并注意由第一个方程，右端非负。所以 $b'$ 的解总是存在，且至多相差 $\tilde\lambda$ 的作用\footnote{$\tilde\lambda$ 在 $b'$ 上的作用除点 $b'=u=v=0$ 外是自由的，在该点 $\tilde\lambda$ 的作用与 $\lambda$ 生成同一条轨道。所以在 $d=1/2$ 时，把 $b'$ 固定为非负实数后，就可以从描述中略去 $\tilde\lambda$。}。于是 $d=1/2$ 处的辛商由 $u,v,b$ 参数化，并受 $|u|^2+|v|^2+2|b|^2=1$ 以及 $|\lambda|=1$ 时 $(u,v,b)\cong (\lambda u,\lambda v,\lambda^2 b)$ 的等价关系约束。合起来就给出加权射影空间 $\Bbb{WCP}^2(1,1,2)$。

反过来做，把 $d$ 从 $1/2$ 减小到更小的值，就解出了 $\Bbb{WCP}^2(1,1,2)$ 的奇点，得到 Hirzebruch 曲面。此外，这张特定的 Hirzebruch 曲面可以形变为简单的乘积 $\Bbb{CP}^1\times \Bbb{CP}^1$。为此只需注意：丛 $\CO(2)\oplus \CO$ 不稳定，可以形变为 $\CO(1)\oplus \CO(1)$。但 $\CO(1)\oplus \CO(1)$ 的射影化与 $\CO\oplus \CO$ 的射影化相同，因而是 $\Bbb{CP}^1$ 上的平凡 $\Bbb{CP}^1$ 丛，即乘积 $\Bbb{CP}^1\times \Bbb{CP}^1$。

结论是：与一般的奇异代数簇不同，$\Bbb{WCP}^2(1,1,2)$ 的奇点在形变或解奇后给出拓扑上相同的结果，而且形变与解奇可以同时进行。此外，这也与奇点附近的局部结构是超凯勒 $A_1$ 奇点这一事实有关。在第~\ref{sec:monbub}~节末尾以及第~\ref{sec:kahstr}~节中，我们将从 Bogomolny 方程的观点解释这一行为。我们还将在第~\ref{sec:opprod}~节利用这个奇点附近的结构来展示 't Hooft 算子的算子乘积展开的计算。

\begin{yizhu}[物理图像]
按：$\Bbb{WCP}^2(1,1,2)$ 的奇点局部是超凯勒 $A_1$ 奇点 $\Bbb{C}^2/\Bbb{Z}_2$。超凯勒结构使形变（$a_1a_3-a_2^2=\epsilon b^2$，即把奇点换成 $\Bbb{CP}^1\times\Bbb{CP}^1$）与解奇（把辛商参数 $d$ 从 $1/2$ 减小，得到 Hirzebruch 曲面）可以同时进行；对一般的复奇点，这两件事通常不能相容地同时做。物理上，这正是 't Hooft 算子插入点附近的几何可以用 Bogomolny 方程来描述的原因，后文第~\ref{sec:monbub}~节的单极子冒泡图景会再次用到它。
\end{yizhu}

\par\medskip{\it 一般情形：$U(2)$}

型为 $^L\negthinspace w=(m_1,m_2)$ 的 Hecke 修正空间可以类似地分析。首先，型 $(m_1+c,m_2+c)$ 的 Hecke 修正空间与型 $(m_1,m_2)$ 的相同，因为加上 $(c,c)$ 只是把 Hecke 变换的输出与 ${\cal O}(p_0)^c$ 作张量积。所以只需取 $^L\negthinspace w=(a,0)$，其中 $a$ 是某个整数，并可假定非负。（否则先把 $(a,0)$ 换成 $(0,a)$，再给每个权加上 $-a$，就换成 $(-a,0)$。）在描述型 $(a,0)$ 的 Hecke 修正时，对于在 $z=0$ 处有 $n$ 阶极点的函数 $g$，方便的做法是用 $[g]_-$ 记 $g$ 的极部：$[g]_-=g_{-n}z^{-n}+g_{-n+1}z^{-n+1}+\dots + g_{-1}z^{-1}$。

相对于上面引入的分解 $E_-={\cal L}_1\oplus {\cal L}_2$，用 $^L\negthinspace w=(a,0)$ 作 Hecke 变换，得到丛 $E_+$，它的一般截面取形式
% 式 (9.23)
\setcounter{equation}{22}
\begin{equation}\label{eq:indon}
s=[c_0s_1]_-+[c_1zs_1]_-+[c_2z^2s_1]_-+\dots
+[c_{a-1}z^{a-1}s_1]_-+(h_1(z),h_2(z)),
\end{equation}
其中复系数为 $c_0,c_1,\dots,c_{a-1}$。

因此 $E_+$ 只依赖于 $s_1$ 模去 $z^a$ 阶的项。这留下总共 $2a$ 个复参数。然而，若把 $s_1$ 换成 $gs_1$（其中 $g$ 是在 $z=0$ 处可逆的全纯函数），丛 $E_+$ 不受影响。通过选取 $g$，可以消去 $s_1$ 中 $a$ 个复参数，这是对 $a=1,2$ 时做法的推广。所以型 $(a,0)$ 的 Hecke 修正空间复维数为 $a$。为了紧化这个空间，必须把型 $(a-1,1), \,(a-2,2),\dots, (a-[a/2],[a/2])$ 的 Hecke 修正也包括进来，它们出现在 $s_1$ 的 Taylor 级数展开中相继的项为零之时。

若把 $^L\negthinspace w=(a,0)$ 换成 $^L\negthinspace w=(a/2,-a/2)$，我们已解释的内容中有些对多数物理学家会更清楚。（Hecke 修正空间在这个平移下当然不变。）正如我们在第~\ref{sec:theck}~节 所解释的，权 $(a/2,-a/2)$ 对规范群 $G=PSU(2)=SO(3)$、$^L\negthinspace G=SU(2)$ 才有意义，而且具有这些权的 't Hooft 算子与 $SU(2)$ 的自旋 $a$ 表示中的 Wilson 圈互为 $S$-对偶。

在研究 $SU(2)$ 规范理论中的禁闭时，考虑基于任意表示 $R$ 的 Wilson 圈是多余的。对中心作用相同的两个表示而言，在这个用途下是等价的，所以 $SU(2)$ 唯一需要考虑的非平凡表示是自旋 $1/2$ 表示。对偶地，$SO(3)$ 的任何 't Hooft 算子都在拓扑上等价于平凡算子或具有权 $(1/2,-1/2)$ 的算子。正是因为更高的 't Hooft 算子在拓扑上不稳定，在旨在理解规范理论相的那些 't Hooft 算子应用（例如 't Hooft 的原始工作）中，它们被认为无关紧要。

从我们的观点看，不稳定的 't Hooft 算子很重要，但它们可以与同一拓扑型的更低算子「混合」。在第~\ref{sec:monbub}~节 中，我们将用 Bogomolny 方程解释这种混合意味着什么。

这种混合的一般性质在 Lie 群 $G$ 与 $^L\negthinspace G$ 为单群时最容易表述。设 $^L\negthinspace R$ 是 $^L\negthinspace G$ 的一个表示。若 $^L\negthinspace G$ 的表示 $^L\negthinspace R'$ 的最高权是 $^L\negthinspace R$ 的权但不是最高权，我们就说 $^L\negthinspace R'$ 与 $^L\negthinspace R$ 相伴。记 $^L\negthinspace w$ 与 $^L\negthinspace w'$ 分别为 $^L\negthinspace R$ 与 $^L\negthinspace R'$ 的最高权，在这种情形下我们也说 $^L w'$ 与 $^L\negthinspace w$ 相伴。条件是：$^L\negthinspace w-{}^L\negthinspace w'$ 是 $^L\negthinspace G$ 的伴随型的非零支配权。

\begin{yizhu}[易错点]
按：「相伴」只看权，不涉及子表示或函子性：定义要求 $^L\negthinspace R'$ 的最高权是 $^L\negthinspace R$ 的一个权，且不能等于 $^L\negthinspace R$ 的最高权。对单群，这等价于 $^L\negthinspace w-{}^L\negthinspace w'$ 是伴随型的非零支配权。极小表示之所以没有相伴表示，是因为它的最高权是它唯一的支配权，从而不可能再给出别的表示；这解释了为什么极小表示对应的 Hecke 修正空间光滑且紧。
\end{yizhu}

于是一般地，型 $^L\negthinspace R$ 的 Hecke 修正空间 ${{\EUY}}(^L\negthinspace R)$ 的自然紧化包含所有与 $^L\negthinspace R$ 相伴的表示 $^L\negthinspace R'$ 所对应的 Hecke 修正。对 $G=U(N)$（或 $SU(N)$ 或 $PSU(N)$），这句话只是说：对权 $^L\negthinspace w=(m_1,m_2,\dots,m_N)$，若对某对 $i,j$ 有 $m_i\geq m_j+2$，则自然紧化包含把 $m_i$ 换成 $m_i-1$、把 $m_j$ 换成 $m_j+1$ 所得的 Hecke 修正，以及通过重复这一移动能得到的所有其他权。这个移动对紧化的必要性，可以简单地在 $U(N)$ 中嵌入一族 $U(2)$ Hecke 修正而看出，在那里权 $(m_i,m_j)$ 退化为 $(m_i-1,m_j+1)$。对更一般的群，可以用 $^L\negthinspace G$ 的各种 $SU(2)$ 子群生成的移动作类似的论证。这些移动在紧化中的作用也可以从仿射 Grassmann 流形看出（我们在第~\ref{sec:affgrass}~节 简要引入这个概念），或者如第~\ref{sec:monbub}~节 那样从 Bogomolny 方程看出。

\par\medskip{\it 极小表示}

把 $U(2)$ 的权 $(1,0)$ 作推广，我们再描述几个 Hecke 修正空间光滑且紧的例子。对 $G=U(N)$，如 \eqref{eq:nyto}\ 那样考虑权 $^L\negthinspace w(k)=(1,1,\dots,1,0,0,\dots,0)$，其中 1 的个数为 $k$。取分解 $E_-=\oplus_{i=1}^N{\cal L}_i$，其中 ${\cal L}_i$ 由截面 $s_i$ 生成，则 Hecke 算子 $T({}^L\negthinspace w(k))$ 把 $E_-$ 映到 $E_+=\oplus_{i\leq k}\left({\cal L}_i\otimes \CO(p_0)\right)\oplus\left(\oplus_{i>k}{\cal L}_i\right)$。$E_+$ 的一般截面取形式 % 式 (9.24)
\setcounter{equation}{23}
\begin{equation}\label{eq:nson}
s={1\over
z}(a_1,\dots,a_N)+(h_1,\dots,h_N)
\end{equation}
在 $p_0$ 附近，其中 $\vec a=(a_1,\dots,a_N)\in \Bbb{C}^N$ 是 $s_1(0),\dots,s_k(0)$ 的线性组合。因此 $E_+$ 只通过 $s_1(0),\dots,s_k(0)$ 生成的子空间 $W\subset \Bbb{C}^N$ 依赖于 $s_i$ 的选取。$W$ 可以是 $\Bbb{C}^N$ 的任意 $k$ 维子空间，所以型 $^L\negthinspace w(k)$ 的 Hecke 修正空间 ${{\EUY}}(^L\negthinspace w(k))$ 就是 $\Bbb{C}^N$ 中复 $k$-平面的 Grassmann 流形 ${\rm Gr}(k,N)$。

于是显然，型 $^L\negthinspace w(k)$ 的 Hecke 修正空间是光滑且紧的。这是对 $^L\negthinspace w=(1,0)$ 的 $U(2)$ 例子的推广。正如我们在关于式 \eqref{eq:nyto} 时所指出的，$^L\negthinspace w(k)$ 是 $^L\negthinspace G=U(N)$ 的表示 $\wedge^k\VV$ 的最高权。这个表示是「极小」（minuscule）表示，意思是它的权构成单一的 Weyl 轨道。一般地，极小表示对应于稳定的 't Hooft 算子，其 Hecke 修正空间光滑且紧。极小表示恰好是那些没有相伴表示的表示，因为极小表示的最高权是它唯一的支配权。

$SU(N)$ 的类似表示 $\wedge^k\VV$ 的最高权 $^L\negthinspace w'(k)$，是把 $^L\negthinspace w(k)$ 的每个权减去常数 $-k/N$ 而得的。Hecke 修正空间不变。表示 $\wedge^k\VV$ 都是极小的，并且生成 $SU(N)$ 的表示环。然而，$SU(N)$ 是唯一一个表示环由极小表示生成的单连通单 Lie 群。对一般的单群 $G$，$G$ 的中心的每个特征恰对应一个极小表示；$SU(N)$ 之外的群中心相对较小，极小表示也很少。

\begin{yizhu}[原文疑误]
按：原文写「subtracting the constant $-k/N$」（减去常数 $-k/N$），但把每个权减去常数 $-k/N$ 实际上是加上 $k/N$。$\wedge^k\VV$ 作为 $SU(N)$ 的表示，其最高权应由 $^L\negthinspace w(k)$ 减去 $k/N$（使权零和）得到；已发表的原文也写作 $-k/N$，译文照录，读者宜按「减去 $k/N$」理解。
\end{yizhu}

\par\medskip{\it Hecke 修正空间的维数}

对 $U(N)$（或 $SU(N)$，或 $PU(N)$），我们来确定任意型的 Hecke 修正空间的维数。对权 $^L\negthinspace w=(a_1,a_2,\dots,a_N)$，其中 $a_1\geq a_2\geq \dots\geq a_N$，维数为
% 式 (9.25)
\setcounter{equation}{24}
\begin{equation}\label{eq:ydro}
\Delta_{^L\negthinspace w}=\sum_{i<j}(a_i-a_j).
\end{equation}
为论证这一断言，考虑在 $z=0$ 附近生成平凡秩 $N$ 丛 $E_-=\CO\oplus \CO\oplus\dots\oplus \CO$ 的 $N$ 个截面 $s_i$。我们回忆，这意味着 $s_i(0)$ 线性无关。然后考虑一般截面为下式的丛 $E_+$
% 式 (9.26)
\setcounter{equation}{25}
\begin{equation}\label{eq:ytgo}
q=\sum_{i=1}^Nz^{-a_i}g_is_i,
\end{equation}
其中函数 $g_i(z)$ 在 $z=0$ 处全纯。丛 $E_+$ 在下述变换下不变 % 式 (9.27)
\setcounter{equation}{26}
\begin{equation}\label{eq:gunjo}
s_i\to s_i+\sum_{j\leq
i}h_js_j+\sum_{j>i}z^{a_i-a_j}h_js_j,
\end{equation}
其中 $h_j$ 是 $z=0$ 处全纯的一般函数。对 $s_i$ 的这种变换可以吸收进 \eqref{eq:ytgo} 中出现的函数 $g_i$ 的重定义。把 $s_i$ 在 $z=0$ 附近展开成 Taylor 级数，不能被变换 \eqref{eq:gunjo}\ 消去的系数个数为 $\sum_{j=i+1}^N(a_i-a_j)$。对 $i$ 求和，就得到 Hecke 修正空间维数的公式 \eqref{eq:ydro}。
